\documentclass[11pt, a4paper]{article}

% Gemeinsame Style-Definitionen (TEXINPUTS muss templates/ enthalten — siehe scripts/build_tex.py)
% bfw-style.tex — gemeinsame Style-Definitionen für alle Handouts/Aufgabenhefte
% Voraussetzung: Kompilation mit xelatex (wegen fontspec)

\usepackage[ngerman]{babel}
\usepackage{geometry}
\geometry{a4paper, top=2.2cm, bottom=2.2cm, left=2.3cm, right=2.3cm}

\usepackage{fontspec}
% Fallback-Kette: Source Sans Pro -> Calibri -> Segoe UI -> Latin Modern Sans
% (Latin Modern ist immer mit MiKTeX vorhanden)
\IfFontExistsTF{Source Sans Pro}{%
  \setmainfont{Source Sans Pro}[Ligatures=TeX]%
}{\IfFontExistsTF{Calibri}{%
  \setmainfont{Calibri}[Ligatures=TeX]%
}{\IfFontExistsTF{Segoe UI}{%
  \setmainfont{Segoe UI}[Ligatures=TeX]%
}{%
  \setmainfont{Latin Modern Sans}[Ligatures=TeX]%
}}}
\IfFontExistsTF{Source Code Pro}{%
  \setmonofont{Source Code Pro}[Scale=0.92]%
}{\IfFontExistsTF{Consolas}{%
  \setmonofont{Consolas}[Scale=0.92]%
}{%
  \setmonofont{Latin Modern Mono}[Scale=0.92]%
}}

\usepackage{xcolor}
% Farbpalette: hell, freundlich, modern
\definecolor{bfwPrimary}{HTML}{0EA5A4}    % Petrol/Türkis - Primär
\definecolor{bfwPrimaryDark}{HTML}{0F766E} % dunkler Petrol für Headings
\definecolor{bfwAccent}{HTML}{F59E0B}     % Amber - Akzent
\definecolor{bfwSuccess}{HTML}{10B981}    % Grün - Erfolg / Lösung
\definecolor{bfwWarn}{HTML}{F97316}       % Orange - Achtung
\definecolor{bfwInk}{HTML}{0F172A}        % fast Schwarz
\definecolor{bfwInkSoft}{HTML}{475569}    % gedämpftes Grau
\definecolor{bfwBgSoft}{HTML}{F8FAFC}     % off-white
\definecolor{bfwBgTeal}{HTML}{ECFEFF}     % helles Türkis
\definecolor{bfwBgAmber}{HTML}{FEF3C7}    % helles Gelb
\definecolor{bfwBgRose}{HTML}{FFE4E6}     % helles Rosé für Eselsbrücken

\usepackage[colorlinks=true, linkcolor=bfwPrimaryDark, urlcolor=bfwPrimaryDark]{hyperref}
\usepackage{enumitem}
\setlist[itemize]{leftmargin=*, itemsep=2pt, topsep=2pt}
\setlist[enumerate]{leftmargin=*, itemsep=2pt, topsep=2pt}

\usepackage[most]{tcolorbox}
\usepackage{booktabs}
\usepackage{tikz}
\usetikzlibrary{decorations.pathmorphing, positioning, shapes.geometric, arrows.meta}

% Merkkästchen — wichtige Definition / Kernpunkt
\newtcolorbox{merksatz}[1][]{%
  enhanced, breakable,
  colback=bfwBgTeal,
  colframe=bfwPrimary,
  boxrule=0.6pt,
  arc=4pt,
  left=10pt, right=10pt, top=8pt, bottom=8pt,
  fonttitle=\bfseries\color{bfwPrimaryDark},
  title={\faLightbulb\ Merksatz},
  #1
}

% Eselsbrücke — Mnemoniks
\newtcolorbox{eselsbruecke}[1][]{%
  enhanced, breakable,
  colback=bfwBgRose,
  colframe=bfwAccent,
  boxrule=0.6pt,
  arc=4pt,
  left=10pt, right=10pt, top=8pt, bottom=8pt,
  fonttitle=\bfseries\color{bfwWarn},
  title={Eselsbrücke},
  #1
}

% Beispiel-Box
\newtcolorbox{beispiel}[1][]{%
  enhanced, breakable,
  colback=bfwBgSoft,
  colframe=bfwInkSoft,
  boxrule=0.4pt,
  arc=3pt,
  left=10pt, right=10pt, top=8pt, bottom=8pt,
  fonttitle=\bfseries\color{bfwInk},
  title={Beispiel},
  #1
}

% Lösungs-Box (für Aufgabenhefte)
\newtcolorbox{loesung}[1][]{%
  enhanced, breakable,
  colback=bfwBgAmber!40,
  colframe=bfwSuccess,
  boxrule=0.6pt,
  arc=3pt,
  left=10pt, right=10pt, top=8pt, bottom=8pt,
  fonttitle=\bfseries\color{bfwSuccess!70!black},
  title={Lösung},
  #1
}

% Glossar-Eintrag
\newcommand{\glossareintrag}[2]{%
  \par\noindent\textbf{\color{bfwPrimaryDark}#1} \enspace #2\par\smallskip
}

% Section-Stil — etwas farbiger als Standard
\usepackage{titlesec}
\titleformat{\section}
  {\Large\bfseries\color{bfwPrimaryDark}}
  {\thesection}{0.6em}{}
\titleformat{\subsection}
  {\large\bfseries\color{bfwPrimary}}
  {\thesubsection}{0.5em}{}
\titleformat{\subsubsection}
  {\normalsize\bfseries\color{bfwInk}}
  {\thesubsubsection}{0.4em}{}

% TikZ-Stil "skizzenhaft" — etwas weicher als Rough.js, aber stabil
\tikzset{
  roughbox/.style = {
    draw=bfwPrimaryDark, thick, fill=bfwBgTeal,
    rounded corners=4pt, inner sep=6pt
  },
  roughaccent/.style = {
    draw=bfwAccent, thick, fill=bfwBgAmber,
    rounded corners=4pt, inner sep=6pt
  },
  roughline/.style = {
    draw=bfwInkSoft, thick, -{Stealth[length=2.5mm]}
  }
}

% Bullet-Symbole (bewusst kein FontAwesome — vermeidet Font-Installations-Probleme)
\newcommand{\faLightbulb}{$\bullet$}
\newcommand{\faBookmark}{$\bullet$}

% Header/Footer
\usepackage{fancyhdr}
\pagestyle{fancy}
\fancyhf{}
\renewcommand{\headrulewidth}{0pt}
\fancyfoot[L]{\small\color{bfwInkSoft}\thepage}
\fancyfoot[R]{\small\color{bfwInkSoft} BFW M-064 Beschaffung im Büromanagement}


\title{\vspace*{-1.5cm}\Huge\bfseries\color{bfwPrimaryDark}Tag~3: Arbeitsumgebung \& Umweltfaktoren im Büro\\[0.4em]
  \large\color{bfwInkSoft}\textnormal{Klima, Licht, Lärm, Farbe \& Co. --- wie die Umgebung die Leistung prägt}}
\author{\textsf{Büroprozesse gestalten und Arbeitsvorgänge organisieren \textperiodcentered{} M-067}\\
\textsf{\small\copyright\ Can Siebert\,\textperiodcentered\,2026}}
\date{01.07.2026}

\begin{document}
\maketitle
\thispagestyle{empty}

\vspace{1em}
\begin{merksatz}[title={\bfseries Worum geht es heute?}]
Auch das beste Büro nützt wenig, wenn es zu warm, zu laut oder zu dunkel ist. Zur
\emph{Arbeitsumgebung} gehören die Faktoren Klima, Beleuchtung, Lärm, Farbe, Pflanzen,
Strahlung und Gefahrstoffe --- sie alle wirken sich auf Gesundheit, Wohlbefinden und
Leistungsfähigkeit der Beschäftigten aus. In diesem Handout lernen Sie, welche Werte ein
gutes Raumklima auszeichnen, wie viel Licht am Bildschirm nötig ist, wie Lärm gemessen
und gemindert wird und welche Wirkung Farben und Pflanzen haben. Das ist nicht nur
prüfungsrelevant, sondern direkt am eigenen Arbeitsplatz anwendbar.
\end{merksatz}

\tableofcontents
\vspace{1em}
\hrule
\vspace{1em}

\section{Die Arbeitsumgebung im Überblick}

Zur \emph{Arbeitsumgebung} gehören die Faktoren \textbf{Klima, Beleuchtung, Lärm, Farbe,
Pflanzen, Strahlung} und \textbf{Gefahrstoffe}. Alle diese Komponenten wirken sich auf
die Gesundheit und das Wohlbefinden der Mitarbeiter aus und können deren
Leistungsfähigkeit beeinflussen. Wer Büroarbeit organisiert, muss diese Faktoren kennen
und gezielt steuern --- denn schlechte Umgebungsbedingungen kosten Konzentration,
Motivation und am Ende bares Geld.

\begin{merksatz}
Die Arbeitsumgebung ist kein „weicher” Faktor: Schlechtes Klima, Lärm oder Blendung
führen messbar zu Fehlern, Kopfschmerzen, höherem Krankenstand und sinkender
Produktivität. Gute Bedingungen sind daher zugleich Gesundheitsschutz und
Wirtschaftlichkeit.
\end{merksatz}

\section{Raumklima}

Das Klima setzt sich u.\,a.\ aus drei Bestandteilen zusammen: \textbf{Lufttemperatur},
\textbf{relative Luftfeuchtigkeit} und \textbf{Luftgeschwindigkeit} (Luftbewegung). Die
Arbeitsstättenverordnung (ArbStättV) schreibt für einen Arbeitsplatz ein gesundheitlich
zuträgliches Raumklima vor. Konkretisiert wird dies durch die Technischen Regeln für
Arbeitsstätten (ASR); die \emph{ASR A3.5} beschäftigt sich explizit mit der
\textbf{Raumtemperatur}.

\subsection{Lufttemperatur}

In Büroräumen soll eine \emph{Mindesttemperatur von +20\,\textdegree C} vorherrschen.
Arbeitsräume \textbf{sollen} +26\,\textdegree C nicht überschreiten; geschieht dies doch
(z.\,B.\ bei Sonneneinstrahlung), sind zunächst Jalousien einzusetzen, ansonsten
zusätzliche Maßnahmen wie Gleitzeitregelungen zur Arbeitszeitverlegung, Lockerung der
Bekleidungsregeln oder das Bereitstellen geeigneter Getränke (z.\,B.\ Wasser). Werden
+30\,\textdegree C überschritten, \textbf{müssen} zusätzliche Maßnahmen erfolgen.

\begin{eselsbruecke}
\textbf{„20 -- 26 -- 30: soll, soll, muss.”} Bei \emph{20}\,\textdegree C beginnt das
zuträgliche Klima, ab \emph{26}\,\textdegree C \emph{soll} gehandelt werden, ab
\emph{30}\,\textdegree C \emph{muss} gehandelt werden. Das Temperaturempfinden ist dabei
individuell --- es hängt z.\,B.\ von Geschlecht, Gesundheitszustand, Bekleidung,
Gewöhnung und Tätigkeit ab.
\end{eselsbruecke}

\subsection{Luftfeuchtigkeit und Luftbewegung}

Die \textbf{Luftfeuchtigkeit} gibt an, wie hoch der Wasserdampfgehalt der Luft ist, und
wird mit dem \emph{Hygrometer} gemessen. Eine relative Luftfeuchtigkeit von
\textbf{40 -- 60\,\%} gilt als optimal. Ist sie zu hoch (insbesondere im Winter), droht
\emph{Schimmelgefahr}; ist sie zu niedrig, führt das zu brennenden Augen sowie trockenen
Lippen und Schleimhäuten --- wodurch die Gefahr der Ansteckung durch
Infektionskrankheiten ansteigt.

Die \textbf{Luftbewegung} wird vor allem durch Klimaanlagen, schlecht isolierte Fenster
oder ungünstig angebrachte Lüftungsschächte verursacht. Eine Luftbewegung von
\emph{0,1 m/s bis 0,15 m/s} wird von den meisten Menschen noch als angenehm empfunden.
Herrscht dagegen \emph{Zugluft} vor, kann dies zu Nacken- und Rückenschmerzen führen.

\subsection{Klimatisierung und Lüften}

Von \emph{Klimatisierung} darf nur gesprochen werden, wenn \textbf{alle drei Faktoren
automatisch geregelt} werden. Klimaanlagen verbrauchen sehr viel Strom und können bei
schlechter Wartung Krankheitserreger übertragen.

\begin{merksatz}
Um das Raumklima zu beeinflussen, sollten in nicht klimatisierten Räumen die Fenster
\textbf{stündlich zum Stoßlüften} geöffnet werden; ansonsten bietet sich die Pausenzeit
zum Lüften an. Weiterhin kann das Raumklima durch \emph{Heizen, Kühlen} oder
\emph{Entlüften} beeinflusst werden.
\end{merksatz}

\section{Licht und Arbeitsplatzbeleuchtung}

Da nicht an jedem Arbeitsplatz genügend Tageslicht vorhanden ist --- bzw.\ das
vorhandene Tageslicht je nach Jahres- und Uhrzeit nicht ausreicht --- muss auf
\emph{künstliche Beleuchtung} zurückgegriffen werden. Beleuchtung beeinflusst das
Wohlbefinden und die Arbeitsleistung. Ist sie ergonomisch durchdacht, lassen sich
Kopfschmerzen, Augenbeschwerden, Nervosität und Ermüdung vermeiden. Wichtig sind dabei:

\begin{itemize}
  \item Die Beleuchtung muss an das \textbf{Sehvermögen} des Nutzers und an die jeweilige
    Tätigkeit \emph{angepasst} werden können.
  \item Der Arbeitsraum sollte \textbf{gleichmäßig hell} sein, damit sich die Augen nicht
    ständig neu einstellen müssen (\,$=$ Anstrengung).
  \item Die \textbf{Lichtfarbe} (Farbtemperatur) bestimmt, wie das Licht wirkt. Für
    Bildschirmarbeit wird eine \emph{neutral weiße} oder \emph{tageslichtweiße} Lichtfarbe
    empfohlen.
  \item \textbf{Reflexblendungen} (z.\,B.\ vom Bildschirm) und \textbf{Direktblendungen}
    (z.\,B.\ von Lampen im Blickfeld) sollen vermieden werden. Hilfreich ist es, den
    Bildschirm \emph{parallel zum Fenster} aufzustellen.
  \item Für Bildschirmarbeitsplätze, Arbeits- und Besprechungsräume wird eine
    Beleuchtungsstärke von \textbf{mindestens 500 Lux} gefordert (Lux $=$ Einheit für die
    Beleuchtungsstärke), ergänzt durch persönliche Schreibtischleuchten.
  \item Die Beleuchtung muss \textbf{flimmerfrei} sein.
\end{itemize}

\begin{beispiel}
Ein Sachbearbeiter klagt über Kopfschmerzen und tränende Augen. Der Bildschirm steht mit
dem Rücken zum Fenster, sodass sich die Fensterfläche darin spiegelt (Reflexblendung).
Lösung: Bildschirm um 90\,\textdegree\ drehen (parallel zum Fenster), Sonnenschutz
(Jalousie) gegen direkte Einstrahlung, blendfreie Deckenleuchte mit mindestens 500 Lux
und eine zusätzliche Schreibtischleuchte für das Lesen auf Papier.
\end{beispiel}

\section{Lärm}

Gespräche, Geräusche und Töne werden als \emph{Lärm} empfunden, wenn sie jemanden stören
und belästigen. Das führt zu einer sehr individuellen Bewertung: Wer ein Lied mag,
empfindet es als Genuss --- ein anderer als Lärm. Auch der persönliche
Gesundheitszustand beeinflusst die Lärmempfindlichkeit.

Lärm wird in \textbf{Dezibel --- dB(A)} --- gemessen; dies gibt an, wie hoch der
sogenannte \emph{Schalldruckpegel} ist. Sind Menschen dauerhaft Lärm ausgesetzt, kann
dies zu erheblichen Gesundheitsbeeinträchtigungen führen: Konzentration und
Leistungsfähigkeit leiden, und der Stressfaktor steigt. Deshalb befasst sich auch die
ArbStättV in ihrem Anhang~3.7 mit Lärm: Der Schalldruckpegel ist in Abhängigkeit von der
Art des Betriebes so niedrig wie möglich zu halten. Ab einem Wert von
\textbf{85 dB(A)} muss den Arbeitnehmern ein geeigneter, persönlicher Gehörschutz zur
Verfügung gestellt und auch von ihnen getragen werden.

\begin{merksatz}
Die Bundesanstalt für Arbeitsschutz und Arbeitsmedizin (BAuA) bewertet den
Schalldruckpegel am Bildschirmarbeitsplatz bis \textbf{30 dB(A)} als \emph{optimal}, bis
\textbf{40 dB(A)} als \emph{sehr gut} und bis \textbf{45 dB(A)} als \emph{gut}
(„Arbeitswissenschaftliche Erkenntnisse Nr.~124: Lärmminderung in Mehrpersonenbüros”).
\end{merksatz}

\subsection{Maßnahmen zur Lärmreduktion}

\begin{itemize}
  \item Bei der Auswahl von Geräten (Drucker, Kopierer) auf \emph{geräuscharme Typen}
    achten (Herstellerangaben beachten).
  \item Laute Geräte in \emph{separaten Räumen} unterbringen.
  \item \emph{Schallschutzhauben} für Geräte anschaffen.
  \item Möglichkeiten zur \emph{Schalldämpfung} nutzen: Teppiche, Vorhänge oder
    schallabsorbierende Deckenplatten.
\end{itemize}

\begin{beispiel}
\textbf{Modern im Großraum:} Wo offene Flächen viele Menschen zusammenbringen, helfen
heute zusätzlich \emph{Akustikkabinen} (Telefon-/Fokuszellen) für Telefonate und
konzentriertes Arbeiten, \emph{Schallabsorber-Stellwände} zwischen den Arbeitsplätzen
sowie \emph{Akustikdecken}. So sinkt der Geräuschpegel, ohne die Vorteile der offenen
Fläche aufzugeben.
\end{beispiel}

\section{Farbgestaltung}

Je nachdem, wie ein Arbeitsplatz oder -raum farblich gestaltet ist, erzeugt dies
unterschiedliche psychologische Wirkungen. Farben wirken auf Stimmung, Leistung und
Wohlbefinden und sollten aufeinander abgestimmt werden. Grundsätzlich gilt:

\begin{itemize}
  \item \textbf{Warme Farben} (Rot, Orange, Gelb) wirken \emph{aktivierend und
    stimulierend} und lassen einen Raum kleiner wirken.
  \item \textbf{Kalte Farben} (Blau, Grün, Violett) wirken \emph{beruhigend und
    distanzierend} und lassen einen Raum größer wirken.
  \item Je kräftiger die Farbe, desto stärker ihre Wirkung.
\end{itemize}

Eine sinnvolle Verwendung von Farbe in Büroräumen kann u.\,a.\ die \emph{Motivation}
steigern, die Wahrnehmung ungünstiger Arbeitsbedingungen positiv beeinflussen, die
\emph{Energiekosten} senken (da das Temperaturempfinden beeinflusst wird), die
Sicherheits- und Gesundheitsschutzkennzeichnung unterstützen und die Erholung fördern.
Wichtig: Farbe und Beleuchtung dürfen nicht unabhängig voneinander betrachtet werden, da
sie sich gegenseitig in ihrer Wirkung beeinflussen. Oftmals gibt auch das
\emph{Corporate Design} bestimmte Farben vor (z.\,B.\ „Telekom Magenta”).

\subsection{Farben und Sicherheitszeichen}

Farben spielen auch bei den \emph{Sicherheitszeichen} eine große Rolle, da sie dem
Gesundheitsschutz und der Unfallverhütung dienen:

\begin{center}
\begin{tabular}{p{3.0cm} p{4.0cm} p{4.7cm}}
\toprule
\textbf{Farbe} & \textbf{Art des Zeichens} & \textbf{Bedeutung / Beispiel}\\
\midrule
Rot & Verbotszeichen & Untersagt Handlungen (z.\,B.\ „Rauchen verboten”)\\
Rot & Brandschutzzeichen & Kennzeichnet Feuermelde- und Feuerbekämpfungseinrichtungen\\
Gelb & Warnzeichen & Weist auf Gefahren hin (z.\,B.\ explosionsgefährliche Stoffe)\\
Blau & Gebotszeichen & Schreibt Verhalten vor (z.\,B.\ „Augenschutz benutzen”)\\
Grün & Rettungszeichen & Markiert Rettungswege, Erste Hilfe, Sammelpunkte\\
\bottomrule
\end{tabular}
\end{center}

\begin{eselsbruecke}
\textbf{„Rot verbietet, Gelb warnt, Blau gebietet, Grün rettet.”} So merken Sie sich die
vier Grundfarben der Sicherheitskennzeichnung.
\end{eselsbruecke}

\section{Pflanzen, Strahlung und Gefahrstoffe}

\subsection{Pflanzen (Raumbegrünung)}

Pflanzen im Büro sorgen allein durch ihr Dasein für \emph{Entspannung und Wohlbefinden}.
Gleichzeitig reinigen sie die Luft von Schadstoffen (je nach Pflanzenart z.\,B.\
Formaldehyd, Benzol, Kohlendioxid) und Chemikalien (wie Lacke oder Reinigungsmittel). Sie
produzieren Sauerstoff, erhöhen die Luftfeuchtigkeit und \emph{schlucken Schall und
Staub}. Dadurch kann die Gesundheit positiv beeinflusst werden: Der Lärmpegel sinkt,
Atemwegserkrankungen wird vorgebeugt, und Erkältungsviren breiten sich langsamer aus.

\begin{beispiel}
\textbf{Biophiles Design:} Moderne Bürokonzepte holen die Natur bewusst ins Gebäude ---
etwa mit begrünten Wänden (Moos- oder Pflanzenwänden), die zugleich als
\emph{Schallabsorber} wirken, das Raumklima verbessern und das Wohlbefinden steigern.
\end{beispiel}

\subsection{Strahlung}

Elektromagnetische Strahlung entsteht in der Umgebung von elektrischen Leitungen und
Geräten (Computer, Smartphones, DECT-Telefone, W-LAN, Bluetooth). Ihr werden
Gesundheitsrisiken nachgesagt; die Forschung ist hier jedoch noch nicht zu einem
endgültigen Ergebnis gelangt. Die ArbStättV verlangt, dass die von Bildschirmgeräten
ausgehende Strahlung so niedrig gehalten wird, dass die Gesundheit der Beschäftigten
nicht gefährdet wird. Maßnahmen zur Reduktion:

\begin{itemize}
  \item \emph{Strahlungsarme Geräte} verwenden (Flachbildschirme statt Röhrenbildschirme).
  \item Im Büro auf tragbare Telefone verzichten und feste Telefone verwenden.
  \item Nicht gebrauchte Geräte abschalten.
  \item W-LAN-Access-Points mindestens \emph{3\,m} vom Arbeitsplatz entfernt anbringen.
\end{itemize}

\subsection{Gefahrstoff im Büro: Tonerstaub}

Auch in Büroräumen können Gefahrstoffe auftreten. Insbesondere ist dies
\textbf{Tonerstaub}, der aus Laserdruckern, Laserfax- und -kopiergeräten in die Luft
gelangt --- etwa beim Wechseln der Kartuschen, bei Wartungsarbeiten oder beim Beseitigen
eines Papierstaus. Über die gesundheitlichen Auswirkungen weiß man noch wenig; mögliche
Folgen sind schwere Atemwegserkrankungen, Hautprobleme, Reizung der Schleimhäute,
Reizhusten, Dauerschnupfen oder Augenbrennen. Manche Toner enthalten zudem krebserregende
Stoffe. Vorsichtsmaßnahmen:

\begin{itemize}
  \item Bei der Anschaffung auf das \emph{DGUV-Testzeichen} für besonders emissionsarme
    Geräte achten und nur \emph{schadstoffarme Toner} mit Prüfsiegel kaufen.
  \item Geräte in separaten, gut belüfteten Räumen aufstellen.
  \item Verschütteten Toner sofort \emph{feucht} wegwischen (nicht aufwirbeln).
  \item Nach dem Beseitigen eines Papierstaus gründlich die Hände waschen.
  \item Regelmäßige Wartung durch Fachpersonal; Tonerkartuschen nicht gewaltsam öffnen;
    Geräte ggf.\ mit Feinstaubfiltern nachrüsten.
\end{itemize}

\section{Zusammenfassung}

\begin{enumerate}
  \item Zur \textbf{Arbeitsumgebung} gehören Klima, Beleuchtung, Lärm, Farbe, Pflanzen, Strahlung und Gefahrstoffe --- alle beeinflussen Gesundheit und Leistung.
  \item \textbf{Raumklima}: Mindesttemperatur 20\,\textdegree C, ab 26\,\textdegree C soll, ab 30\,\textdegree C muss gehandelt werden; Luftfeuchtigkeit 40--60\,\% optimal; stündlich stoßlüften.
  \item \textbf{Beleuchtung}: mindestens 500 Lux, neutral-/tageslichtweiße Lichtfarbe, blend- und flimmerfrei, Bildschirm parallel zum Fenster.
  \item \textbf{Lärm}: gemessen in dB(A); ab 85 dB(A) Gehörschutz; am Bildschirm optimal bis 30, sehr gut bis 40, gut bis 45 dB(A); Maßnahmen: geräuscharme/separate Geräte, Schallschutz.
  \item \textbf{Farbe}: warme Farben aktivieren, kalte beruhigen; Sicherheitszeichen --- Rot verbietet, Gelb warnt, Blau gebietet, Grün rettet.
  \item \textbf{Pflanzen} verbessern Luft, Akustik und Wohlbefinden (biophiles Design); \textbf{Strahlung} und \textbf{Tonerstaub} durch geeignete Maßnahmen minimieren.
\end{enumerate}

\newpage

\section*{Glossar}
\addcontentsline{toc}{section}{Glossar}

\glossareintrag{Arbeitsumgebung}{Gesamtheit der Faktoren Klima, Beleuchtung, Lärm, Farbe, Pflanzen, Strahlung und Gefahrstoffe, die auf Gesundheit und Leistung wirken.}

\glossareintrag{ArbStättV}{Arbeitsstättenverordnung; schreibt u.\,a.\ ein gesundheitlich zuträgliches Raumklima und Schutz vor Lärm und Strahlung vor.}

\glossareintrag{ASR A3.5}{Technische Regel für Arbeitsstätten, die die Raumtemperatur konkretisiert (Mindesttemperatur, Grenzwerte 26/30\,\textdegree C).}

\glossareintrag{Raumklima}{Zusammenspiel aus Lufttemperatur, relativer Luftfeuchtigkeit und Luftgeschwindigkeit.}

\glossareintrag{Hygrometer}{Messgerät für die relative Luftfeuchtigkeit (optimal 40--60\,\%).}

\glossareintrag{Stoßlüften}{Kurzzeitiges, kräftiges Lüften bei weit geöffnetem Fenster, möglichst stündlich, zur Verbesserung des Raumklimas.}

\glossareintrag{Klimatisierung}{Automatische Regelung aller drei Klimafaktoren; nur dann darf von Klimatisierung gesprochen werden.}

\glossareintrag{Lux}{Einheit für die Beleuchtungsstärke; am Bildschirmarbeitsplatz mindestens 500 Lux gefordert.}

\glossareintrag{Lichtfarbe}{Auch Farbtemperatur; bestimmt, wie Licht wirkt --- für Bildschirmarbeit neutral-/tageslichtweiß empfohlen.}

\glossareintrag{Blendung}{Sehbeeinträchtigung durch Reflexblendung (z.\,B.\ am Bildschirm) oder Direktblendung (Lampe im Blickfeld).}

\glossareintrag{Dezibel dB(A)}{Maßeinheit für den Schalldruckpegel; gibt an, wie laut ein Geräusch ist.}

\glossareintrag{Schalldruckpegel}{Maß für die Lautstärke; am Bildschirmarbeitsplatz optimal bis 30, sehr gut bis 40, gut bis 45 dB(A).}

\glossareintrag{Gehörschutz}{Persönliche Schutzausrüstung; ab 85 dB(A) bereitzustellen und zu tragen.}

\glossareintrag{Farbpsychologie}{Lehre von der psychologischen Wirkung von Farben (warm $=$ aktivierend, kalt $=$ beruhigend).}

\glossareintrag{Sicherheitszeichen}{Farbcodierte Zeichen: Rot $=$ Verbot/Brandschutz, Gelb $=$ Warnung, Blau $=$ Gebot, Grün $=$ Rettung.}

\glossareintrag{Raumbegrünung}{Einsatz von Pflanzen, die Luft reinigen, Sauerstoff produzieren, Schall und Staub schlucken und das Wohlbefinden steigern.}

\glossareintrag{Biophiles Design}{Gestaltungsansatz, der Natur (z.\,B.\ Pflanzenwände) ins Gebäude holt und Klima, Akustik und Wohlbefinden verbessert.}

\glossareintrag{Tonerstaub}{Gefahrstoff aus Laserdruckern/-kopierern; kann Atemwege und Schleimhäute reizen, manche Toner sind krebserregend.}

\end{document}
